\chapter{概形的态射}\label{src:A}
\begin{epigraph}
概形论通常被看作一种极其抽象的代数工具, 它遮蔽了几何的魅力, 以追求令人难以招架且往往没有必要的一般性. 相反, 专家知道, 概形使事情变得更简单.
\epigraphauthor{David Eisenbud 与 Joe Harris \cite{EH00}}
\end{epigraph}
我们回顾概形态射的定义与性质: 局部有限表现, 平坦, 光滑, 平展, 非分歧, 拟有限和固有. 为突出模理论中的重要结果, 我们特别关注函子性质, 即由概形及其态射的函子刻画的性质.

尽管代数几何的入门课程 (例如 \cite{Har77}) 常常施加 Noether 假设, 我们仍尽可能在非 Noether 情形下陈述结果. 由于模函子和叠定义在整个概形范畴上, 使用非 Noether 概形是必不可少的. 极限方法 (\bookref{B.3}) 使我们能够把一般概形的性质归结为 Noether 概形的性质.
\setlength{\emergencystretch}{4em}
\hypersetup{pageanchor=false}
\AddToHook{shipout/foreground}[appendix-a-index-pages]{\put(0,0){\makebox[0pt][l]{\hypertarget{page.\thepage}{}}}}
